% ============================================================
% 3. THIRD QUESTION
% ============================================================
\section{Regression Analysis}\label{sec:q3}

% ------------------------------------------------------------
% a)
% ------------------------------------------------------------
\subsection{Best Simple Linear Regression Model}\label{sec:q3a}

\task{Create a data frame considering only the variables: body fat, grip force, sit ups counts, gender and age categories.
Consider the numerical variables in the data set and find the best simple linear regression model to predict the broad jump variable.
Test the assumptions and use transformations if it is required.
Explain why the model you find is the best simple linear regression model and interpret the model coefficients.}

\paragraph{Data.}
We created the data frame \texttt{q3} with the response \texttt{broad\_jump\_cm} and the five variables of the task (\texttt{body\_fat\_pct}, \texttt{grip\_force}, \texttt{sit\_ups}, \texttt{gender} and \texttt{age\_cat}), from the 12,742 observations of the cleaned data set (Section~\ref{sec:q1b}) and with the age groups of Section~\ref{sec:q1e}.

\paragraph{Best simple model.}
Among the three numerical variables, \texttt{grip\_force} has the strongest correlation with \texttt{broad\_jump\_cm} ($r = 0.751$), closely followed by \texttt{sit\_ups} ($r = 0.743$), both strong and positive, while \texttt{body\_fat\_pct} has a moderate negative correlation ($r = -0.677$) (\autoref{tab:q3a-correlation}).
\texttt{grip\_force} gives the model with the highest $R^2$ and the smallest residual standard error of the three (\autoref{tab:q3a-slr} and \autoref{fig:q3a-scatter}).
It explains 56.4\% of the variability of \texttt{broad\_jump\_cm}, with a typical error of 25.73\,cm.
Therefore, \texttt{grip\_force} is the numerical variable that gives the best simple linear regression model, \texttt{broad\_jump\_cm} $\sim$ \texttt{grip\_force}.

\paragraph{Assumptions.}
We first checked the assumptions of the simple linear regression on this model without any transformation (\autoref{tab:q3a-assumptions}, row \texttt{Original}):
\begin{enumerate}
  \item \textbf{Linearity:} the scatter plot shows a linear trend between \texttt{grip\_force} and \texttt{broad\_jump\_cm} (\autoref{fig:q3a-scatter}).
  \item \textbf{Normality of the error term:} the Shapiro-Wilk test on a random sample of 5,000 residuals rejects normality at a 95\% confidence level ($p = 2.19 \times 10^{-13} < 0.05$), but the Q-Q plot follows the line except in the tails (\autoref{fig:q3a-diagnostics}, right), and with this sample size the test detects even very small deviations.
  \item \textbf{Homoscedasticity:} the Breusch-Pagan test rejects homoscedasticity at a 95\% confidence level ($p = 9.41 \times 10^{-6} < 0.05$; residuals against fitted values in \autoref{fig:q3a-diagnostics}, left).
  \item \textbf{Independence of the error term:} the Durbin-Watson test does not reject the absence of autocorrelation at a 95\% confidence level ($p = 0.468 > 0.05$).
\end{enumerate}

\paragraph{Box-Cox transformation.}
Since homoscedasticity does not hold, we applied the Box-Cox transformation to \texttt{broad\_jump\_cm} ($\lambda = 1.46$, \autoref{fig:q3a-boxcox}).
With the transformed response, the Breusch-Pagan test still rejects homoscedasticity at a 95\% confidence level ($p = 2.67 \times 10^{-5} < 0.05$), and the other assumptions do not change (\autoref{tab:q3a-assumptions}, row \texttt{Box-Cox}).
Since the transformation does not solve the heteroscedasticity, we kept the model without transformation.
In its residual plot the spread looks constant (\autoref{fig:q3a-diagnostics}, left), so the heteroscedasticity detected by the test is small.
The 633 outliers (5.0\%) are about the 5\% expected under normality, so we kept them.

\paragraph{Interpretation.}
The fitted model (\autoref{tab:q3a-slr}) is
\[
  \texttt{broad\_jump\_cm} = 87.97 + 2.777\,\texttt{grip\_force}.
\]
The t-test of the slope shows that \texttt{grip\_force} has a significant positive linear effect on \texttt{broad\_jump\_cm} at a 95\% confidence level (R output, Appendix~\ref{app:code}).
1\,kg more of \texttt{grip\_force} increases the predicted \texttt{broad\_jump\_cm} by 2.78\,cm.
The intercept has no practical meaning, since the smallest \texttt{grip\_force} in the data is 4.4\,kg (\autoref{tab:q1b-summary-clean}).

% ------------------------------------------------------------
% b)
% ------------------------------------------------------------
\subsection{Multiple Linear Regression Model}\label{sec:q3b}

\task{Fit a multiple linear regression model with the most important two (numerical) variables.
Use transformations if it is needed and test all the assumptions.
Then compare this model to the simple linear regression model that you fit in (a).
Which one is a better model?
Why?}

\paragraph{Multiple linear regression model.}
Among the three models with two numerical predictors, the model with \texttt{grip\_force} and \texttt{sit\_ups} has the highest adjusted $R^2$ (\autoref{tab:q3b-selection}).
Therefore, \texttt{grip\_force} and \texttt{sit\_ups} are the two most important numerical variables, and the fitted model is
\[
  \texttt{broad\_jump\_cm} = 72.64 + 1.790\,\texttt{grip\_force} + 1.297\,\texttt{sit\_ups}.
\]
The t-tests show that both predictors have a significant positive linear effect on \texttt{broad\_jump\_cm} at a 95\% confidence level (R output, Appendix~\ref{app:code}).
Keeping \texttt{sit\_ups} constant, 1\,kg more of \texttt{grip\_force} increases the predicted \texttt{broad\_jump\_cm} by 1.79\,cm, and keeping \texttt{grip\_force} constant, one more sit-up increases it by 1.30\,cm.

\paragraph{Assumptions.}
We checked the assumptions of the multiple linear regression on this model without any transformation (\autoref{tab:q3b-assumptions}):
\begin{enumerate}
  \item \textbf{Linearity:} the scatter plots show a linear trend between \texttt{broad\_jump\_cm} and each predictor (\autoref{fig:q3b-pairs}), and the residuals stay around 0 for all the fitted values, with only a slight bend at the highest fitted values, above 250\,cm (\autoref{fig:q3b-diagnostics}, left).
  \item \textbf{Normality of the error term:} the Shapiro-Wilk test rejects normality at a 95\% confidence level ($p < 10^{-16} < 0.05$), but the Q-Q plot follows the line except in the tails (\autoref{fig:q3b-diagnostics}, right).
  \item \textbf{Homoscedasticity:} the Breusch-Pagan test does not reject homoscedasticity at a 95\% confidence level ($p = 0.517 > 0.05$), and the residuals form a band of constant width around 0 (\autoref{fig:q3b-diagnostics}, left).
  \item \textbf{Independence of the error term:} the Durbin-Watson test does not reject the absence of autocorrelation at a 95\% confidence level ($p = 0.802 > 0.05$).
  \item \textbf{No multicollinearity:} \texttt{grip\_force} and \texttt{sit\_ups} are only moderately correlated ($r = 0.573$, \autoref{fig:q3b-pairs}), and their VIF is 1.49, far below 10.
\end{enumerate}
Since all the assumptions hold except normality, we did not transform the response.
The 614 outliers (4.8\%) are about the 5\% expected under normality, so we kept them.

\paragraph{Comparison with (a).}
We compared this model with the simple model of (a), both without transformation.
Adding \texttt{sit\_ups} increases the adjusted $R^2$ from 0.564 to 0.710 and reduces the residual standard error from 25.73 to 20.99\,cm, and the partial F-test shows that this improvement is significant at a 95\% confidence level ($p < 10^{-16} < 0.05$, \autoref{tab:q3b-comparison}).
In \autoref{fig:q3b-vs-a} the points of (b) are closer to the line $y = x$.
Therefore the multiple model of (b) is better: it explains 71.0\% of the variability of \texttt{broad\_jump\_cm}, against 56.4\% in (a), with smaller errors and without any transformation.

% ------------------------------------------------------------
% c)
% ------------------------------------------------------------
\subsection{Models with a Factor}\label{sec:q3c}

\task{Now we will add a factor to the regression model you have chosen in section (b).
(i) First, add age categories (age\_cat) as a factor.
Interpret the model coefficients and overall summary of the model.
(ii) Add gender to the model that you had chosen in the previous step in (b).
Interpret the model coefficients and overall summary of this model similar to the previous case.}


\subsubsection{Model with Age Categories}\label{sec:q3c-age}

The fitted model with \texttt{age\_cat} (reference category \texttt{<30}, \autoref{tab:q3c-age-coef}) is
\begin{align*}
  \texttt{broad\_jump\_cm} = {}& 83.48 + 1.859\,\texttt{grip\_force} + 1.063\,\texttt{sit\_ups} \\
  & - 1.965\,\texttt{age\_cat30-39} - 4.052\,\texttt{age\_cat40-49} \\
  & - 13.399\,\texttt{age\_cat50+}.
\end{align*}

All the coefficients are significant at a 95\% confidence level.
Keeping \texttt{grip\_force} and \texttt{sit\_ups} constant, people aged 30--39 jump 1.97\,cm less than those under 30, people aged 40--49 4.05\,cm less, and people aged 50 or more 13.40\,cm less.
So the jump distance decreases with age, much more from the age of 50 (\autoref{fig:q3c-categories}, top).
Keeping \texttt{sit\_ups} and \texttt{age\_cat} constant, 1\,kg more of \texttt{grip\_force} increases the predicted \texttt{broad\_jump\_cm} by 1.86\,cm, and keeping \texttt{grip\_force} and \texttt{age\_cat} constant, one more sit-up increases it by 1.06\,cm.


The model is significant at a 95\% confidence level (overall F-test, $p < 10^{-16} < 0.05$; R output, Appendix~\ref{app:code}).
It explains 72.3\% of the variability of \texttt{broad\_jump\_cm}, with a residual standard error of 20.52\,cm (\autoref{tab:q3c-summary}).

\subsubsection{Model with Gender}\label{sec:q3c-gender}

The fitted model with \texttt{gender} (reference category \texttt{Female}, \autoref{tab:q3c-gender-coef}) is
\begin{align*}
  \texttt{broad\_jump\_cm} = {}& 89.76 + 0.943\,\texttt{grip\_force} + 1.278\,\texttt{sit\_ups} \\
  & + 23.608\,\texttt{genderMale}.
\end{align*}

All the coefficients are significant at a 95\% confidence level.
Keeping \texttt{grip\_force} and \texttt{sit\_ups} constant, men jump 23.61\,cm further than women.
So men perform better than women in the broad jump (\autoref{fig:q3c-categories}, bottom).
Keeping \texttt{sit\_ups} and \texttt{gender} constant, 1\,kg more of \texttt{grip\_force} increases the predicted \texttt{broad\_jump\_cm} by 0.94\,cm, and keeping \texttt{grip\_force} and \texttt{gender} constant, one more sit-up increases it by 1.28\,cm.

The model is significant at a 95\% confidence level (overall F-test, $p < 10^{-16} < 0.05$; R output, Appendix~\ref{app:code}).
It explains 74.0\% of the variability of \texttt{broad\_jump\_cm}, with a residual standard error of 19.87\,cm (\autoref{tab:q3c-summary}).

% ------------------------------------------------------------
% d)
% ------------------------------------------------------------
\subsection{Model Comparison}\label{sec:q3d}

\task{Compare the model in section (b) with the two models that had constructed with an additional factor in (c).
Which one would you choose?
Why?}

\paragraph{Comparison.}
We compared the model of (b) with the two models of (c) by the partial F-test of each factor against the model of (b), the residual standard error and the adjusted $R^2$, which takes into account that the models have different numbers of coefficients (\autoref{tab:q3d-comparison}).
Both factors improve the model of (b) significantly at a 95\% confidence level (partial F-tests, $p < 10^{-16} < 0.05$).
Between adding \texttt{gender} or \texttt{age\_cat}, \texttt{gender} gives the largest improvement, with only one coefficient more instead of the three of \texttt{age\_cat}.

Visually, the difference between the three models is hardly noticeable in \autoref{fig:q3d-fit}, but it is clear in their adjusted $R^2$ and residual standard errors in \autoref{tab:q3d-comparison}.

\csvtable[\scriptsize]{q3d_comparison}
  {d) Comparison of the model of b) with the two models of c); partial F-test of each factor against the model of b).}
  {tab:q3d-comparison}

\paragraph{Choice.}
We choose the model of (b) with \texttt{gender} (Section~\ref{sec:q3c-gender}) since it has the highest adjusted $R^2$ and the smallest error, and it improves the model of (b) significantly with only one coefficient more.

% ------------------------------------------------------------
% e)
% ------------------------------------------------------------
\subsection{Validity of the Final Model}\label{sec:q3e}

\task{Test the validity of the final model that you choose.}

\paragraph{Assumptions.}
We checked the assumptions of the final model, the model of (b) with \texttt{gender} (Section~\ref{sec:q3c-gender}), on its residuals (\autoref{tab:q3e-validity} and \autoref{fig:q3e-residuals}):
\begin{enumerate}
  \item \textbf{Linearity:} the residuals stay around 0 for all the fitted values (\autoref{fig:q3e-residuals}, top left).
  \item \textbf{Normality of the error term:} the Shapiro-Wilk test rejects normality at a 95\% confidence level ($p < 10^{-16} < 0.05$), but the histogram and the Q-Q plot are close to normal except in the tails, especially the lower one, which gives the histogram a slight left skew (\autoref{fig:q3e-residuals}, top right and bottom left).
  \item \textbf{Homoscedasticity:} the Breusch-Pagan test rejects homoscedasticity at a 95\% confidence level ($p = 1.11 \times 10^{-14} < 0.05$).
    The Levene test shows that the variance of the residuals is different for women and men ($p = 2.76 \times 10^{-9} < 0.05$): the residuals of men are more spread, with a standard deviation of 20.5\,cm against 18.6\,cm for women (\autoref{fig:q3e-residuals}, bottom right).
  \item \textbf{Independence of the error term:} the Durbin-Watson test does not reject the absence of autocorrelation at a 95\% confidence level ($p = 0.686 > 0.05$).
  \item \textbf{No multicollinearity:} the largest VIF of the predictors is 3.23, below 10.
\end{enumerate}
The 591 outliers (4.6\%) are about the 5\% expected under normality, so we kept them.

\paragraph{Robustness.}
Because of the heteroscedasticity, we also computed the standard errors of the coefficients without assuming homoscedasticity (\autoref{tab:q3e-robust}).
They are only slightly larger than the usual ones, and every coefficient is still significant at a 95\% confidence level ($p < 10^{-16} < 0.05$).

\paragraph{Conclusion.}
We conclude that the final model is valid: the relation is linear, the errors are independent, the predictors are not collinear and the residuals are close to normal.
Despite the heteroscedasticity, the coefficients remain significant at a 95\% confidence level with heteroscedasticity-robust standard errors.

